\subsection*{Inherited dynamics from the conditioned activation law}

\propositionheading{Dynamics of the conditioned parent}{prop:phil-actual-parent-dynamics}
Use the selected matching, channel, and conditioned activation of
\cref{prop:phil-actual-conditioned-parent}. More explicitly, let
\(0<\mathrm T\le1/100\), \(s=t_{\mathrm L}\), \(\theta=s+1\), and let \(\mathscr E\)
be the selected coherence event on \([s,\theta]\) at resolution
\(\mathsf a\) from \cref{prop:phil-selected-pair-persistence}.
Assume \(\Pr(\mathscr E^c)\le\rho<1\), and, throughout
\([s,\theta+\mathsf h]\subseteq[t_{\mathrm L},2t_{\mathrm L}]\),
\[
 \operatorname{TV}(q_t,\nu_{\mathrm T}^M)\le\delta,
 \qquad \max_i|p_i(t)-p_i^M|\le\delta.
\]
Write \(\varepsilon=\epsilon_{\mathrm L}\),
\(d=\delta+\varepsilon\mathsf h<1/4\), and use relative future time
\(u=t-\theta\). Initialize the parent with
\[
 q_0^C=C_\#\operatorname{Law}(\mathbf X_\theta,S_\theta\mid\mathscr E),
 \qquad p^C(0)=p(\theta).
\]
On the stratum \(\{S\}\times W_S\), its continuous generator is
\(-\mathsf A_S y\cdot\nabla_{W_S}+\mathrm T\Delta_{W_S}\).
Its rate-one proposal kernel is
\begin{equation}
 J_p((S,y),\cdot)=
 \int\kappa_S(d\mathbf x\mid y)
     \sum_{S'}b_p(S'\mid\mathbf x)
                   \delta_{(S',P_{W_{S'}}\mathbf x)},
 \label{eq:phil-actual-parent-kernel}
\end{equation}
where \(b_p(S'\mid\mathbf x)=\mu_{\mathrm T,p}(S'\mid\mathbf x)\) is the
product of receiver posterior rows at rows \(p\), and
\(\kappa_S\) is the optimizing hidden conditional: the three
independent Gaussian coordinates on a matching stratum, and a point
mass off the matchings. The parent rows evolve from its own law:
\begin{equation}
 \frac d{du}p_i^C(u)
 =\varepsilon\bigl(q_u^C(\sigma_i=+1)-p_i^C(u)\bigr).
 \label{eq:phil-actual-parent-learning}
\end{equation}
This coupled parent system has a unique global law with interior rows.

Let \(\mathsf P_{\theta,\mathsf h}^{\mathrm{fine},C}\) be the law of
\((C(\mathbf X_{\theta+u},S_{\theta+u}))_{0\le u\le\mathsf h}\)
conditioned only on \(\mathscr E\), with the fine process continuing
its original row path. Let \(\mathsf P_{\mathsf h}^C\) be the parent
path law from the displayed activation and its own learning equation.
Then
\begin{equation}
 \operatorname{TV}
 (\mathsf P_{\theta,\mathsf h}^{\mathrm{fine},C},
                                      \mathsf P_{\mathsf h}^C)
 \le2\rho+(\mathsf h+3)\delta+
                         \frac{4\mathsf h d}{1-4d}.
 \label{eq:phil-actual-parent-path-error}
\end{equation}
The comparison includes all retained coordinates and the assignment
history, hence also the loop history. Self-proposals remain part of
the rate-one kernel.

Along \cref{cor:phil-joint-selection}, choose
\(\delta\le\mathrm T\), \(\varepsilon\le\mathrm T^2\),
\(\mathsf a=\mathrm T^{1/4}\), and
\(\mathsf h=\mathrm T^{-1/4}\). One may take
\(\rho=3\delta+\mathsf b(\mathrm T,\mathsf a,1)\), where
\(\mathsf b\) is from \cref{prop:phil-selected-pair-persistence}.
The right side of \cref{eq:phil-actual-parent-path-error} is then
\(O(\mathrm T^{1/4})\). This is a conditional, finite-window
approximation after the selected coherence event; the activation law
is not replaced by a Gaussian law.
\status{ESTABLISHED}

\noindent\emph{Proof.}
For a matching \(M\), the decomposition
\(\R^{12}=W_M\oplus Z_M\) of
\cref{prop:rg-formation-endogenous-cycle} is preserved by
\(H_M\), \(\Omega\), and \(\mathsf A_M\).
The hidden drift matrix is
\(\operatorname{diag}(8+\mathrm T,4+5\mathrm T,4+5\mathrm T)\).
The prior weight \(\pi_S(p)\) cancels from the conditional law at
fixed \((S,y)\), so \(\kappa_S\) is independent of \(p\).
Off the matchings no coordinate is discarded. Thus the displayed
continuous generator is the projection of every held-sheet diffusion,
and \cref{eq:phil-actual-parent-kernel} is a probability kernel.
A self-proposal has \(S'=S\) and
\(P_{W_S}\mathbf x=y\), so it leaves the parent state unchanged.
At a different assignment it uses that assignment's projection,
including its twelve-dimensional stratum when appropriate.

For fixed rows, the pushed optimizing law \(C_\#\mu_{\mathrm T,p}\)
is invariant. Indeed each diffusion is reversible on its Gaussian
stratum, and
\(\mu_{\mathrm T,p}(d\mathbf x,S)b_p(S'\mid\mathbf x)\)
is symmetric in \(S,S'\). Pushing both endpoints through \(C\)
gives detailed balance for \(J_p\).
The loop covector annihilates each \(Z_M\), so this construction
retains the loop through the stratum changes as well.

We next check the existence assertion with evolving parent rows.
For any continuous interior row path, concatenate the linear
diffusions with the rate-one proposals, drawing \(\kappa_S\) at
each proposal. On every finite interval the clock has finitely many
rings and every inserted Gaussian is finite almost surely, so this
construction is nonexplosive. Its second moments remain finite:
projection cannot increase the norm, an inserted hidden draw adds
expected squared norm at most \(5\mathrm T/8\), and the diffusion
generator applied to \(|y|^2\) is at most
\(24\mathrm T-2\mathrm T|y|^2\). The full generator is consequently
bounded above by \(25\mathrm T-2\mathrm T|y|^2\).
The activation law has finite second moment because it conditions a
finite-moment fine law on an event of positive probability.

If two row paths have entries in \([\mathrm c_0,1-\mathrm c_0]\), couple their
parent processes by the same initial state, Brownian increments,
clocks, hidden Gaussian draws, and receiver uniforms until their
first different proposal. The derivative of a posterior receiver
probability with respect to its prior entry is at most
\(1/[4\mathrm c_0(1-\mathrm c_0)]\). Therefore
\[
 \operatorname{TV}(q_u^{C,p},q_u^{C,\widetilde p})
 \le\frac1{\mathrm c_0(1-\mathrm c_0)}
       \int_0^u\|p(v)-\widetilde p(v)\|_\infty dv.
\]
The variation-of-constants map for
\cref{eq:phil-actual-parent-learning} consequently has Lipschitz
constant at most \(\varepsilon u_*^2/[2\mathrm c_0(1-\mathrm c_0)]\) on an interval
of length \(u_*\). On a short enough interval it is a contraction
on continuous row paths in an interior cube. Its image stays in
that cube after further shortening the interval, since the row
velocity has absolute value at most \(\varepsilon\).
Moreover its solution satisfies
\[
 p_i^C(u)\ge p_i^C(0)e^{-\varepsilon u},\qquad
 1-p_i^C(u)\ge(1-p_i^C(0))e^{-\varepsilon u}.
\]
On any prescribed finite horizon these lower bounds give a common
interior cube and a positive contraction step length. Successive
steps give existence and uniqueness on that horizon, and then on
all finite horizons. The source marginal is continuous because its
jump rate is bounded, so the resulting rows are continuously
differentiable. This constructs learning from the parent's own law,
without imposing the fine row path after activation.

Let \(\overline\nu_{\mathrm T}^M\) be the retained Gaussian on
\(W_M\). Write \(\mathsf O^C_{\mathrm T,\mathsf h}\) for its
stationary held-\(M\) path law. The projection of
\(\mathsf O_{\mathrm T,\mathsf h}^M\) from
\cref{prop:phil-selected-pair-persistence} is exactly this law.
Only the learned-arm conclusions of that proposition are used, and
their proof does not use its frozen-arm hypothesis.
Conditioning any probability law on an event of probability
\(1-r\) changes its total variation by at most \(r\) after
any measurable projection. Thus the conditioned native retained
future obeys
\begin{equation}
 \operatorname{TV}
 (\mathsf P_{\theta,\mathsf h}^{\mathrm{fine},C},
                     \mathsf O^C_{\mathrm T,\mathsf h})
 \le\rho+(\mathsf h+2)\delta.
 \label{eq:phil-native-conditioned-held}
\end{equation}
This applies to a prefix event through \(\theta\). No condition is
placed on future coherence or future proposals. The same argument
at the activation time gives
\(\operatorname{TV}(q_0^C,
 \overline\nu_{\mathrm T}^M\otimes\delta_M)\le\rho+\delta\).

For the parent's own rows let \(\zeta_i^C\) be the probability in
row \(i\) assigned to the nonpartner of \(M\). Its equation and the
bound on any probability give
\[
 0\le\zeta_i^C(u)\le\zeta_i^C(0)+\varepsilon u
                  \le\delta+\varepsilon u\le d.
\]
The row-independent normalized sheet weights are at most the matching
weight one for \(0<\mathrm T\le1/100\) by
\cref{prop:phil-temperature-tracking}; hence, for any fixed row vector
with this bound,
\[
 \mu_{\mathrm T,p}(M)\ge\pi_M(p)
               \ge1-\sum_i\zeta_i\ge1-4d.
\]
Define the parent departure probability at a proposal by
\(k_p(y)=J_p((M,y),\{S'\ne M\})\).
Detailed balance, and the fact that the optimizing conditional
retained law on \(M\) is \(\overline\nu_{\mathrm T}^M\)
independently of the rows, give
\begin{equation}
 \int k_p(y)\overline\nu_{\mathrm T}^M(dy)
 \le\frac{1-\mu_{\mathrm T,p}(M)}{\mu_{\mathrm T,p}(M)}
 \le\frac{4d}{1-4d}.
 \label{eq:phil-parent-averaged-departure}
\end{equation}
The first inequality bounds stationary flux from \(M\) to its
complement by the complement's stationary mass.

Maximally couple the parent activation with the stationary retained
matching law and use common Brownian increments, clocks, and parent
proposal draws up to the first parent departure. Before departure,
the continuous paths agree because the matching drift splits, and
self-proposals do not move the retained coordinate. The row path
\(p^C(u)\) constructed above is deterministic, even though it was
computed from the parent's own law. Under the stationary comparison
path, its marginal at every time is
\(\overline\nu_{\mathrm T}^M\). Bound the first-departure
compensator with its stopped indicator by the compensator evaluated
on that whole stationary path. Applying
\cref{eq:phil-parent-averaged-departure} at each deterministic
\(p^C(u)\) gives
\[
 \Pr(\text{parent departure before }\mathsf h
                         \text{ while initially coupled})
 \le\frac{4\mathsf h d}{1-4d}.
\]
No Gaussian assertion is made about the parent's law conditioned on
survival or about the native hidden variables. We have proved
\[
 \operatorname{TV}(\mathsf P_{\mathsf h}^C,
                         \mathsf O^C_{\mathrm T,\mathsf h})
 \le\rho+\delta+\frac{4\mathsf h d}{1-4d}.
\]
Combining this with \cref{eq:phil-native-conditioned-held} proves
\cref{eq:phil-actual-parent-path-error}.

For the displayed joint choices,
\(d\le\mathrm T+\mathrm T^{7/4}\le2\mathrm T\),
\(\mathsf h\delta\le\mathrm T^{3/4}\), and
\(\rho=O(\mathrm T^{1/4})\).
The activation and following window fit inside the selection window
because \(1+\mathrm T^{-1/4}\le t_{\mathrm L}\) for sufficiently
small temperature under the mobility cap. The departure term is
\(O(\mathrm T^{3/4})\), giving the stated total error.
The Gaussian comparison laws served only to bound that error; the
parent's activation and subsequent learning remain the laws in the
statement. \hfill\(\square\)


